\section{Limitations and Conclusion}
\label{sec:limits}

\paragraph{Committee, epochs and pool.} The seed's proposer can bias the lottery, slots go to the fastest claimants, and a rushing dealer biases pool keys by choosing which of its dealings enter $\QUAL$ (\S\ref{sec:security}). Finalization has no deadline and is one $2.3$~M-constraint proof ($2.3$~s, $5.5$~GB) that some node must produce before any key of the epoch exists---an epoch-wide liveness dependency on one prover, served in seed-derived stagger slots, where an earlier design went live without a proof; a member's contribution prover peaks at $9.3$~GB for $3.5$~s per epoch, and $\Kmax$ is a release constant that doubled both when doubled. Re-registration protects only future encryptions. An epoch serves at most $\Kmax$ applications, further ones waiting for the next; $\Kmax - 1$ permissionless registrations let $\mathsf{createEpoch}$ run early at the committee's expense, which also pays $t$ partials per admitted ciphertext, with no fee or admission list to deter either.

\paragraph{Organizer key, submissions and trust base.} $\mathsf{sk}_{\mathtt{org}}$ is a single point of failure: losing it makes a locked application undecryptable, reusing it lets one reveal open every application it backs, and leaking it permits an unauthorized but visible reveal, or a quiet decryption if $t$ members also collude (\S\ref{sec:security}). Admitted submitters can maul same-application ciphertexts, so admission is the application's responsibility; identifiers being first come, first served, clients must verify registrant and key on chain. An out-of-range plaintext costs each combiner a failed discrete-logarithm search; tainting the offending application--submitter pair bounds this per submitter, not against fresh addresses. The single-party setup's toxic waste permits false plaintexts (combine) or a chosen $P_j$ (dealing), the proof library is part of the trusted base, and nothing here is post-quantum; future work: admission economics, beacon-based selection \cite{micali1999vrf,drand2024}, proactive refresh \cite{herzberg1995proactive}.

\paragraph{Conclusion.} Proof-carrying dealings give an EVM contract a publicly verifiable DKG with one fixed-size verification per contribution and no dispute phase; keys dealt as a pool and finalized by one proof isolate the applications one committee serves, and an organizer key adds a one-shot, application-wide release.
